\chapter{Error Detection and Correction: CRC and FEC}
\label{ch:error-control}

Every digital signal traveling over a radio link is exposed to noise,
fading, and interference that can flip received bits. This chapter
covers the two complementary families of techniques used to deal with
that reality: detecting that an error occurred, and, where possible,
correcting it without needing a retransmission.

\begin{figure}[H]
\centering
\begin{tikzpicture}[
  node distance=6mm,
  box/.style={draw,rounded corners,minimum width=3.2cm,minimum height=8mm,align=center,font=\small},
  hdr/.style={font=\bfseries},
]
\node[hdr] (txhdr) {Transmitter};
\node[box,below=of txhdr,fill=white] (tx1) {Original data};
\node[box,below=of tx1,fill=red!12] (tx2) {Apply FEC};
\node[box,below=of tx2,fill=green!12] (tx3) {Calculate CRC};

\node[hdr,right=4.4cm of txhdr] (rxhdr) {Receiver};
\node[box,below=of rxhdr,fill=green!12] (rx1) {Verify CRC};
\node[box,below=of rx1,fill=red!12] (rx2) {Correct with FEC};
\node[box,below=of rx2,fill=white] (rx3) {Received data};

\draw[-{Latex}] (tx1) -- (tx2);
\draw[-{Latex}] (tx2) -- (tx3);
\draw[-{Latex}] (rx1) -- (rx2);
\draw[-{Latex}] (rx2) -- (rx3);
\draw[-{Latex},dashed] (txhdr.east) -- node[above,font=\footnotesize]{radio channel} (rxhdr.west);
\end{tikzpicture}
\caption{Error detection (CRC) vs.\ forward error correction (FEC) in a
digital transmission chain.}
\end{figure}

\section{The Basic Problem}

A block of $k$ information bits, transmitted over a noisy channel, may
arrive with one or more bits flipped. Without any additional structure,
the receiver has no way to know this happened --- an altered but
still-valid-looking bit pattern is simply accepted as correct. Error
control techniques add carefully structured \textbf{redundant} bits to
the transmitted data specifically so the receiver can recognize (and, in
some schemes, correct) errors.

\section{Cyclic Redundancy Check (CRC)}

A \textbf{cyclic redundancy check} is an error-\emph{detection} (not
correction) scheme. The transmitter treats the data block as a large
binary number, divides it (using a special binary polynomial division,
modulo 2) by a fixed \textbf{generator polynomial}, and appends the
remainder --- the \textbf{CRC checksum} --- to the transmitted data. The
receiver performs the same division on the received data (including the
appended checksum) and checks whether the remainder is zero: if not, at
least one bit was altered in transit, and the receiver knows to discard
the block or request retransmission.

\begin{keyidea}
CRC is extremely effective at \emph{detecting} that an error occurred
(essentially guaranteed to catch any single-bit or short burst error,
and to catch the vast majority of longer or multiple-bit error
patterns, depending on the generator polynomial's length), but on its
own gives no way to determine \emph{which} bit was wrong, and therefore
no way to fix it. CRC-protected protocols typically respond to a failed
check by requesting a retransmission (an \textbf{automatic repeat
request}, ARQ) rather than attempting a repair.
\end{keyidea}

CRC is used throughout amateur digital communications: packet radio
(AX.25) frames include a CRC (technically an FCS, frame check sequence)
to catch corrupted packets; many file transfer and digital mode
protocols use CRC in a similar role.

\section{Forward Error Correction (FEC)}

\textbf{Forward error correction} takes a different, more ambitious
approach: it adds redundancy structured so that the receiver can
\emph{correct} at least some error patterns without needing any
retransmission at all --- valuable on links (satellite, weak-signal HF,
one-way beacon or telemetry transmissions) where requesting a repeat is
slow, expensive, or simply impossible.

\subsection{A Simple Example: The Repetition Code}

The simplest possible FEC scheme repeats each bit multiple times (e.g.,
three times: $0 \to 000$, $1 \to 111$) and uses \textbf{majority voting}
at the receiver to decide the original bit, correcting any single-bit
error within each group of three at the cost of tripling the
transmitted data. This illustrates the fundamental FEC trade-off in its
simplest form: \textbf{more redundancy buys more error-correcting
capability, at the direct cost of reduced information throughput for a
given channel bandwidth}.

\subsection{Practical FEC Codes}

Real FEC systems use far more efficient (less redundant, for the same
correcting power) mathematical structures than simple repetition:

\begin{itemize}
\item \textbf{Block codes} (Hamming, Reed--Solomon, BCH): operate on
fixed-size blocks of data, adding a calculated set of parity/check
symbols; Reed--Solomon codes are particularly good at correcting
\textbf{burst errors} (many consecutive bits corrupted together, typical
of fading and impulse noise) and are used in satellite telemetry and
digital broadcast standards.
\item \textbf{Convolutional codes}: encode a continuous stream of data
using a sliding-window structure, decoded with algorithms such as the
\textbf{Viterbi algorithm}; well suited to the continuously varying
noise conditions typical of a real radio channel, and long used in
satellite and deep-space communication.
\item \textbf{Low-density parity-check (LDPC) codes} and modern
\textbf{turbo codes}: iteratively decoded codes that approach the
theoretical (Shannon) limit of channel capacity very closely; used in
recent weak-signal amateur digital modes and modern satellite/deep-space
links.
\end{itemize}

\begin{keyidea}[Why weak-signal digital modes decode below the noise floor]
Modern weak-signal amateur digital modes (Chapter~\ref{ch:digital-modes})
combine strong FEC coding with narrow bandwidth and long averaging
(integration) times to reliably decode signals that are, by ear,
completely inaudible in the noise --- often 10--20\,dB below the noise
floor. The FEC code does not create information out of nothing; rather,
it lets the decoder use redundancy spread across the whole transmitted
block to statistically reconstruct the most likely original message
even when individual symbols are severely corrupted by noise.
\end{keyidea}

\section{Combining CRC and FEC}

Many real protocols use both techniques together: FEC coding to correct
the errors it can, followed by a CRC check on the FEC-decoded output to
catch (and flag or reject) any errors that were too severe for the FEC
to fix. This layered approach --- correct what you can, detect what you
cannot correct so bad data is never silently accepted as good --- is the
standard architecture for robust digital data links, including many
amateur satellite and digital-mode protocols.

\section{Automatic Repeat Request (ARQ)}

Where a return channel exists, error \emph{detection} (CRC) can be
combined with automatic retransmission: the receiver acknowledges
correctly received blocks and requests retransmission of blocks that
fail their CRC check. \textbf{ARQ} protocols (used in amateur packet
radio, Winlink, and many other bidirectional digital modes) trade
throughput and latency (a corrupted block must be resent, taking extra
time) for a very high assurance of eventually delivering error-free
data, in contrast to pure FEC, which trades bandwidth for correcting
errors without any retransmission delay but cannot guarantee a
perfect result on a very poor link.

\begin{quickreview}
\item CRC detects (but does not correct) errors by comparing a
transmitted and recomputed checksum; failed checks typically trigger a
retransmission request (ARQ).
\item FEC adds structured redundancy that lets the receiver correct
certain error patterns without retransmission, at the cost of reduced
raw throughput.
\item Block codes (Reed--Solomon, Hamming), convolutional codes, and
modern LDPC/turbo codes are the main practical FEC families, each suited
to different error patterns (burst vs.\ random) and decoding complexity
budgets.
\item Strong FEC plus narrow bandwidth and long integration time lets
modern digital modes decode signals well below the audible noise floor.
\item CRC and FEC are frequently combined: FEC corrects what it can,
CRC flags what remains uncorrected.
\end{quickreview}

\begin{practiceproblems}
\item Explain the key functional difference between error
\emph{detection} and error \emph{correction}.
\item A simple triple-repetition FEC code receives the group ``101''
for a transmitted single bit. Using majority voting, what was the
original bit most likely to have been, and can you be certain?
\item Explain why Reed--Solomon-style codes are particularly well suited
to correcting burst errors caused by fading, compared to a code
optimized only for isolated single-bit errors.
\item Explain why a protocol might use both FEC and CRC together rather
than relying on FEC alone.
\item Explain, at a conceptual level (no equations required), how
combining strong FEC with a very narrow receive bandwidth and long
averaging time allows some digital modes to decode signals inaudible to
the human ear.
\end{practiceproblems}
